\documentclass[11pt]{article}
\pdfinfoomitdate=1
\pdftrailerid{}
\pdfsuppressptexinfo=-1
\ifdefined\pdfvariable\pdfvariable suppressoptionalinfo -1\fi
\usepackage[margin=0.85in]{geometry}
\usepackage{longtable}
\usepackage{array}
\usepackage{hyperref}
\usepackage[utf8]{inputenc}
\usepackage[T1]{fontenc}
\usepackage{lmodern}
\usepackage{enumitem}

\hypersetup{
  colorlinks=true,
  linkcolor=black,
  urlcolor=blue
}

\setlength{\parindent}{0pt}
\setlength{\parskip}{6pt}
\renewcommand{\arraystretch}{1.2}
\newcolumntype{L}[1]{>{\raggedright\arraybackslash}p{#1}}
\emergencystretch=3em
\Urlmuskip=0mu plus 1mu

\title{\textbf{Source Data, Code, and Reproducibility Guide}\\
\large End-to-end autonomous scientific discovery on a real optical platform}

\author{
\scriptsize Shuxing Yang$^{1,2,3}$, Fujia Chen$^{1,2,3}$, Rui Zhao$^{1,2,3}$, Junyao Wu$^{1,2,3}$,\\
\scriptsize Yize Wang$^{1,2,3}$, Haiyao Luo$^{1,2,3}$, Ning Han$^{4}$, Qiaolu Chen$^{1,2,5}$,\\
\scriptsize Yuze Hu$^{1}$, Wenhao Li$^{1,2,3}$, Mingzhu Li$^{6}$, Hongsheng Chen$^{1,2,3,*}$,\\
\scriptsize and Yihao Yang$^{1,2,3,*}$\\[6pt]
\tiny $^{1}$State Key Laboratory of Extreme Photonics and Instrumentation, College of Information Science and Electronic Engineering,\\
\tiny ZJU-Hangzhou Global Scientific and Technological Innovation Center, The Electromagnetics Academy at Zhejiang University,\\
\tiny Zhejiang University, Hangzhou 310027, China\\
\tiny $^{2}$Key Laboratory of Advanced Micro/Nano Electronic Devices \& Smart Systems of Zhejiang,\\
\tiny Jinhua Institute of Zhejiang University, Zhejiang University, Jinhua 321099, China\\
\tiny $^{3}$Shaoxing Institute of Zhejiang University, Zhejiang University, Shaoxing 312000, China\\
\tiny $^{4}$College of Optical and Electronic Technology, China Jiliang University, Hangzhou 310018, China\\
\tiny $^{5}$Laboratory of Wave Engineering, Ecole Polytechnique Fédérale de Lausanne (EPFL), Lausanne, Switzerland\\
\tiny $^{6}$School of Information and Electrical Engineering, Hangzhou City University, Hangzhou 310015, China\\[6pt]
\scriptsize Correspondence should be addressed to Yihao Yang (\texttt{yangyihao@zju.edu.cn})\\
\scriptsize and Hongsheng Chen (\texttt{hansomchen@zju.edu.cn}).
}

\date{}

\begin{document}
\maketitle

\section*{Purpose}

This guide accompanies the Zenodo source-data archive for the manuscript
\textit{End-to-end autonomous scientific discovery on a real optical platform}.
We provide it to help reviewers and readers locate the data, offline analysis
scripts, generated figures, provenance records, and report-level material
supporting the analyses reported in the paper.

The archive is arranged in five report-based sections, with additional files
for selected main-text data-panel replots and a short qualitative video of the
optical platform. The formal main manuscript and the formal Supplementary
Information are not duplicated here; they are submitted separately through the
journal system and remain the authoritative scientific documents.

The files in this archive support inspection and offline rerunning of analyses
from the included data. Re-acquiring optical measurements requires the original
optical platform and instrument-control environment, and is therefore outside
the scope of this source-data/code deposit.

\section*{Repository Deposit and Journal Upload}

Because this source-data package is larger than a typical journal
Supplementary Information file, we deposit it on Zenodo and cite the record in
the journal submission. The submission route is:

\begin{enumerate}[leftmargin=1.5em]
  \item Upload \path{SM_Source_Materials_CORE_5REPORTS.tar.gz} and its
  checksum file to the Zenodo record associated with the DOI below.
  \item Use the same repository DOI in the manuscript, Supplementary
  Information, this guide, and the Zenodo record.
  \item Cite the DOI in the manuscript Data Availability and Code Availability
  statements.
  \item Submit the formal main manuscript and formal Supplementary Information
  through the journal submission system.
  \item If the journal requests a reviewer-facing PDF for the archive, submit
  this guide PDF as the lightweight source-data/code guide.
\end{enumerate}

Repository DOI: \textbf{10.5281/zenodo.19890402}.\\
Repository URL: \url{https://doi.org/10.5281/zenodo.19890402}.

This archive should be cited as the source-data/code and provenance deposit for
the study, and should be read alongside the formal Supplementary Information.

\section*{Relation to the Main Manuscript and SI}

The main manuscript and Supplementary Information carry the final scientific
description: the claims, methods narrative, figure assembly, panel labels, and
captions. The Zenodo archive has a narrower role. It records the data, scripts,
generated outputs, and report-level materials that allow the analyses to be
inspected and rerun offline.

The folder \path{main_text_data_panels/} provides audit replots for selected
data-derived main-text panels only: Fig. 2d-g, Fig. 3d-f, and Fig. 4g-i. The
PDF/PNG files in that folder are standalone replots made from the included data
and scripts. They are not exact visual copies of the final assembled manuscript
panels. Differences in crop, color scale, font, layout, panel spacing,
normalization display, or selected representative channels/tasks are expected
and do not change the role of the underlying data.

The five report PDFs included under \path{*/report/} are retained as source
reports documenting the analysis history and experiment provenance. They do not
replace the formal manuscript or SI. If report-level wording, file labels, or
figure layouts differ from the final submission documents, the manuscript and
SI should be treated as the formal reference, while this package should be used
for data/code traceability.

The same boundary is also summarized in
\path{SM_SCOPE_AND_RELATION_TO_MANUSCRIPT_SI.md}.

\section*{Package Contents}

\begin{longtable}{@{}L{0.34\linewidth}L{0.58\linewidth}@{}}
\textbf{Path} & \textbf{Contents} \\
\hline
\path{SM_SCOPE_AND_RELATION_TO_MANUSCRIPT_SI.md} & Boundary between this repository archive, the submitted main manuscript, and the submitted SI. \\
\path{ZENODO_DEPOSIT_NOTES.md} & Practical notes for repository deposition and DOI linking. \\
\path{00_context/} & Qualitative visual documentation of the optical platform. \\
\path{01_result2_paper/} & Report 1: initial transmission-matrix reproduction. \\
\path{02_result2_refine_paper/} & Report 2: refined self-referenced transmission-matrix experiment. \\
\path{03_result3_paper/} & Report 3: majorization-order transport experiment. \\
\path{04_result4_paper/} & Report 4: Complex-B bilinear interaction field, XOR showcase, and semantic benchmark. \\
\path{05_result4_refine_paper/} & Report 5: refined Complex-B advantage map and sparse-output analyses. \\
\path{main_text_data_panels/} & Audit replots for selected main-text data panels Fig. 2d-g, Fig. 3d-f, and Fig. 4g-i; not exact final manuscript panels. \\
\path{shared_raw_data/result4_frame_data/} & Full Result4 frame-level CMOS \texttt{.npy} arrays shared by Reports 4 and 5. \\
\path{provenance/} & File tree, inventory, checksums, section map, figure-source map, and raw-data inclusion map. \\
\end{longtable}

Each report section contains four standard subfolders:
\path{report/}, \path{source_data/}, \path{scripts/}, and \path{figures/}.

The video \path{00_context/visual_documentation/video_qiushi.mp4} is included
as qualitative visual documentation of the live Qiushi Engine optical-platform
interface, including instrument-status indicators, optical-power monitoring,
SLM/CMOS display panels, and a witness-camera view of the physical setup. It is
not used as quantitative evidence for any reported result.

The paths listed in this guide are the paths to use for review and offline
reruns. Some source reports were generated during the internal analysis
workflow; if a report refers to an older path label, the corresponding files
should be located through the report-specific folders and provenance files in
this archive.

\section*{Integrity Check}

From the package root, verify the included files:

\begin{verbatim}
cd SM_Source_Materials_CORE_5REPORTS
sha256sum -c provenance/checksums_sha256.txt --quiet
\end{verbatim}

A silent return indicates that the package files match the checksum manifest.

\section*{Software Environment}

The offline analysis scripts use standard scientific Python packages:

\begin{itemize}[leftmargin=1.5em]
  \item Python 3.9 or newer
  \item \texttt{numpy}
  \item \texttt{scipy}
  \item \texttt{matplotlib}
  \item \texttt{scikit-learn}
\end{itemize}

The package does not require the optical hardware-control software for offline
analysis. Hardware-control acquisition scripts are intentionally not included.

Recommended shell setup:

\begin{verbatim}
export MPLBACKEND=Agg
cd SM_Source_Materials_CORE_5REPORTS
\end{verbatim}

\section*{Quick Reproduction Map}

\small
\begin{longtable}{@{}L{0.20\linewidth}L{0.18\linewidth}L{0.31\linewidth}L{0.14\linewidth}@{}}
\textbf{Report section} & \textbf{Data} & \textbf{Main scripts} & \textbf{Generated figures} \\
\hline
\path{01_result2_paper} &
\path{source_data/} &
\path{make_core_figures.py}; \path{make_core_figures_v2.py} &
\path{figures/} \\
\path{02_result2_refine_paper} &
\path{source_data/} &
\path{generate_paper_figures.py}; \path{fix_tm_filtering.py} &
\path{figures/} \\
\path{03_result3_paper} &
\path{source_data/} &
\path{generate_v3_figures.py}; \path{generate_raw_data_figure.py}; analysis scripts &
\path{figures/} \\
\path{04_result4_paper} &
\path{source_data/}; \path{shared_raw_data/} &
XOR, semantic, and compressed-budget scripts &
\path{figures/} \\
\path{05_result4_refine_paper} &
\path{source_data/}; \path{shared_raw_data/} &
phase-2, three-bin, advantage-boundary, and sparse-tap scripts &
\path{figures/} \\
\end{longtable}
\normalsize

For exact data/code/figure mapping, see
\path{provenance/FIGURE_SOURCE_MAP.tsv}.

For main-text data-panel mapping, see
\path{main_text_data_panels/MAIN_TEXT_DATA_PANEL_SOURCE_MAP.tsv}. This map
covers only Fig. 2d-g, Fig. 3d-f, and Fig. 4g-i and includes a visual-status
column marking these outputs as audit replots. It does not attempt to
reconstruct the full manuscript page layouts, final assembled panels, or
schematic panels.

\section*{Report-by-Report Reproduction Commands}

The commands below are intended for offline reruns. They overwrite or refresh
outputs inside the corresponding report section. For a non-destructive review,
copy the package to a working directory before running scripts.

\subsection*{Main-Text Data Panels}

\begin{verbatim}
cd SM_Source_Materials_CORE_5REPORTS
python main_text_data_panels/scripts/generate_main_text_data_panels.py
\end{verbatim}

This command regenerates standalone audit-replot outputs for Fig. 2d-g,
Fig. 3d-f, and Fig. 4g-i under
\path{main_text_data_panels/generated_panels/}. These files support data/code
inspection and should not be used as exact visual copies of the final
manuscript panels.

\subsection*{Report 1: Initial Transmission-Matrix Reproduction}

\begin{verbatim}
cd SM_Source_Materials_CORE_5REPORTS/01_result2_paper
python scripts/make_core_figures.py
python scripts/make_core_figures_v2.py
\end{verbatim}

Inputs are in \texttt{source\_data/}. Figures are written to \texttt{figures/}.

\subsection*{Report 2: Refined Self-Referenced Transmission Matrix}

\begin{verbatim}
cd SM_Source_Materials_CORE_5REPORTS/02_result2_refine_paper
python scripts/generate_paper_figures.py
\end{verbatim}

The script uses \path{source_data/tm_256},
\path{source_data/focusing}, \path{source_data/forward_model},
and \path{source_data/ring_screening}. The optional script
\path{fix_tm_filtering.py} recomputes filtered TM diagnostics and can
overwrite filtered arrays.

\subsection*{Report 3: Majorization-Order Transport}

\begin{verbatim}
cd SM_Source_Materials_CORE_5REPORTS/03_result3_paper
python scripts/generate_v3_figures.py
python scripts/generate_raw_data_figure.py
python scripts/analyze_majorization_spectra.py
python scripts/analyze_majorization_spectra_v2.py
\end{verbatim}

Inputs are in \path{source_data/exp_transport_v1},
\path{source_data/exp_transport_v2}, and
\path{source_data/exp_transport_v3}.

\subsection*{Report 4: Complex-B Bilinear Interaction Field}

The full frame-level CMOS arrays are included once under
\path{shared_raw_data/result4_frame_data/}. Commands:

\begin{verbatim}
cd SM_Source_Materials_CORE_5REPORTS/04_result4_paper
RAW=../shared_raw_data/result4_frame_data

python scripts/analyze_complex_B_xor_showcase.py \
  --data-dir $RAW/exp_bilinear_001/exp_bilinear_001 \
  --output-dir source_data/processed/complex_B_xor_showcase_results

python scripts/analyze_complex_B_semantic.py \
  --data-dir $RAW/exp_complex_B_semantic \
  --output-dir source_data/processed/complex_B_semantic_results

python scripts/generate_semantic_bin10_figures.py
python scripts/generate_semantic_bin10_figures_v2.py

python scripts/analyze_semantic_compressed_budget.py \
  --data-dir $RAW/exp_complex_B_semantic \
  --output-dir source_data/processed/semantic_compressed_budget_results
\end{verbatim}

\subsection*{Report 5: Refined Complex-B Advantage Map}

\begin{verbatim}
cd SM_Source_Materials_CORE_5REPORTS/05_result4_refine_paper
python scripts/analyze_refinement_phase2_fast.py
python scripts/regenerate_xor_complex_scatter_bin10.py
python scripts/analyze_three_bin_sparse.py
python scripts/analyze_advantage_boundary.py
python scripts/analyze_bin_anchored_twin.py
python scripts/analyze_sparse_tap_twin.py
\end{verbatim}

These scripts read the shared raw Result4 frame data and write to the local
\path{source_data/processed/} and \path{figures/} folders.

\section*{Expected Reproduction Level}

\begin{longtable}{p{0.30\linewidth}p{0.62\linewidth}}
\textbf{Level} & \textbf{Supported by this package} \\
\hline
Report inspection & Yes. Five report PDFs are included in the five report sections. \\
Main manuscript and SI inspection & Not included in this repository archive. Submit the formal manuscript and SI separately through the journal system. \\
Figure inspection & Yes. Generated figures are included in each report section. \\
Main-text data-panel rerun & Yes. Fig. 2d-g, Fig. 3d-f, and Fig. 4g-i are regenerated by \path{main_text_data_panels/scripts/generate_main_text_data_panels.py}. \\
Physical-platform visual documentation & Yes. A short MP4 recording is included under \path{00_context/visual_documentation/}. \\
Processed-data inspection & Yes. Each report section contains its own \path{source_data/}. \\
Script inspection & Yes. Offline Python scripts are included and syntax-checked. \\
Offline figure/result rerun & Yes, subject to standard Python dependencies and runtime. \\
Result4 raw-frame reprocessing & Yes. Frame-level CMOS arrays are included in \path{shared_raw_data/}. \\
Optical hardware re-acquisition & No. This requires the original physical platform and hardware-control stack. \\
\end{longtable}

\section*{Validation Summary}

At package finalization, the repository archive contained:

\begin{itemize}[leftmargin=1.5em]
  \item Five report-specific sections.
  \item One main-text data-panel reproduction section for Fig. 2d-g, Fig. 3d-f, and Fig. 4g-i.
  \item No main-manuscript or Supplementary Information PDFs inside the repository archive.
  \item One qualitative optical-platform visual-documentation video.
  \item Full Result4 frame-level raw data under \path{shared_raw_data/}.
  \item Twenty-nine offline Python scripts with zero syntax errors.
  \item A checksum manifest in \path{provenance/checksums_sha256.txt}.
\end{itemize}

\section*{Contact}

Correspondence and requests for materials should be addressed to:

\begin{itemize}[leftmargin=1.5em]
  \item Yihao Yang, \texttt{yangyihao@zju.edu.cn}
  \item Hongsheng Chen, \texttt{hansomchen@zju.edu.cn}
\end{itemize}

\end{document}
